\chapter{Related Work}
\label{chapter:related_work}

Research on urban population simulation spans related areas that differ in spatial scale, population representation, and modeling objectives. Crowd simulation primarily addresses movement and interaction in shared physical spaces, from individual pedestrians to aggregated crowd flows~\cite{thalmann2007crowd,ijaz2015hybridSurvey}. Population mobility research instead examines movement across larger geographic and temporal scales and the regularities that characterize travel between locations~\cite{brockmann2006scaling,gonzalez2008understanding,natureAlessandretti2017multiScaleSpatioTemporal}. Multilevel and urban simulation approaches connect aspects of these perspectives by combining representations or supporting populations interacting with structured urban environments~\cite{serana2023multilevelmodeling,balmer2009matsim,taillandier2019gama}.

This chapter reviews these complementary areas to establish the context for the model proposed in this thesis. The discussion progresses from crowd simulation and population mobility to multilevel and urban population simulation, including dynamic urban services and infrastructure. The final section brings these perspectives together to position LODUS with respect to existing approaches.

\section{Crowd Simulation}
\label{related:crowd_simulation}

Crowd simulation studies the movement and interaction of populations in shared physical environments. Models differ principally in how the population is represented, ranging from individual pedestrians to density fields and aggregated groups~\cite{thalmann2007crowd,ijaz2015hybridSurvey}. These representations involve different trade-offs between behavioral detail, computational cost, and population scale, while recent work has further explored behavioral heterogeneity, adaptive representations, and data-driven techniques~\cite{khan2024heterogeneousSurvey,zhang2023collectiveBehavior}.

Microscopic models explicitly represent each member of the crowd and compute its movement individually~\cite{macal2010agentBasedTutorial,thalmann2007crowd}. The Social Force Model proposed by Helbing and Moln\'ar~\cite{helbing1995social}, for example, represents pedestrian motion through forces associated with desired movement, interactions with other pedestrians, and environmental constraints. ORCA~\cite{van2011reciprocal} addresses collision avoidance through reciprocal velocity obstacles, selecting collision-free velocities while distributing avoidance responsibility among interacting agents. BioCrowds~\cite{de2012simulating} instead discretizes the surrounding space into markers that are associated with nearby agents, producing emergent individual territories through which collision-free movement is performed. Although based on different formulations, these approaches maintain the individual pedestrian as the principal unit of simulation.

Microscopic representations also support differences among individuals. Pelechano et al.~\cite{pelechano2005crowd} incorporated psychological states, roles, communication, stress, and decision processes into crowd behavior. Durupinar et al.~\cite{durupinar2015psychological} modeled personality and emotional contagion to distinguish behaviors associated with different crowd types, while Neto et al.~\cite{neto2017giving} combined personality and emotional characteristics with BioCrowds to simulate psychological contagion among agents. Other work investigates physiological, psychological, competitive, and cooperative influences on pedestrian behavior~\cite{narang2015generating,cao2021competitiveCooperative}. Recent reviews consequently identify physical, psychological, social, and behavioral characteristics as important sources of heterogeneity in agent-based crowd simulation~\cite{khan2024heterogeneousSurvey,zhang2023collectiveBehavior}. Movement rules themselves may also adapt to surrounding conditions; Ruprecht et al.~\cite{ruprecht2024adaptiveMovement}, for example, vary the contribution of global navigation and local collision avoidance according to the situation encountered by an agent.

Such models provide detailed control over pedestrian behavior, but the corresponding population state and interactions must be maintained at the individual level. As population size and behavioral complexity increase, this can lead to substantial computational and parameterization requirements. Macroscopic models address this issue by describing collective quantities rather than independently simulating every pedestrian. Hughes~\cite{hughes2002continuum} formulated pedestrian movement as a continuum in which density evolves according to travel costs and destinations, while Treuille et al.~\cite{treuille2006continuum} use dynamic potential fields that account for density and environmental conditions to guide collective flow.

Aggregated representations improve scalability but retain less information about individual population characteristics. BioClouds~\cite{antonitsch2019bioclouds} provides an alternative group-based abstraction in which populations are organized into clouds whose members share properties such as desired density, velocity, and destination. Clouds compete for space according to these collective characteristics, reducing the amount of independently maintained population state. BioClouds is itself a multilevel model and supports transitions to the microscopic BioCrowds representation~\cite{de2012simulating}, illustrating how aggregate and detailed representations can coexist.

Hybrid and multilevel approaches generalize this principle by using different representations according to the needs of the simulation~\cite{ijaz2015hybridSurvey}. Xiong et al.~\cite{xiong2010hybrid} partition the environment into regions simulated with ORCA~\cite{van2011reciprocal} and Continuum Crowds~\cite{treuille2006continuum}, transferring pedestrians between representations through transition areas. Legion~\cite{legion} instead represents populations as groups that can divide when environmental conditions require greater spatial detail and merge when compatible groups can again be represented collectively. Dang et al.~\cite{dang2023multilevelCrowd} similarly investigate a density-based multilevel architecture in which multiple population representations coexist. These approaches demonstrate that population granularity does not need to remain fixed throughout a simulation.

Mesoscopic and data-driven methods provide further alternatives to manually specified microscopic or aggregate dynamics. MPACT~\cite{lemonari2025mpact} derives interpretable mesoscopic behavior profiles from unlabelled trajectory data, mapping spatial observations to higher-level simulation parameters. Li et al.~\cite{Li2025efficienteCrowdSim} use deep reinforcement learning to obtain navigation strategies for crowds in complex environments, whereas Alvarez et al.~\cite{alvarez2025quationfree} use machine learning to bridge microscopic simulations and lower-dimensional representations of macroscopic crowd evolution. These methods modify how behaviors or reduced representations are obtained, although they still operate over either individual agents, aggregate population fields, or intermediate representations and require suitable observations or simulated experience from which the corresponding dynamics can be learned~\cite{lemonari2025mpact,Li2025efficienteCrowdSim,alvarez2025quationfree}.

Overall, crowd simulation offers mature mechanisms for local movement, pedestrian interactions, behavioral heterogeneity, aggregation, and changes in level of detail. Microscopic approaches favor individual detail, whereas macroscopic and group-based models favor population scale; hybrid and learned representations provide different compromises between these objectives~\cite{ijaz2015hybridSurvey,antonitsch2019bioclouds,khan2024heterogeneousSurvey}. Their primary concern, however, remains movement and interaction within physical space. Population dynamics across geographically distributed locations and longer temporal periods are more directly addressed by human mobility research.

\section{Population Mobility}
\label{related:population_mobility}

Population mobility research examines how people move between locations over geographic space and time. Compared with crowd simulation, the emphasis shifts from local pedestrian interactions toward regularities in displacement distance, visitation frequency, destination choice, and movement between meaningful locations~\cite{brockmann2006scaling,gonzalez2008understanding,natureAlessandretti2017multiScaleSpatioTemporal}. Such abstractions support the analysis of populations over neighborhoods, cities, regions, or countries, commonly through statistical descriptions derived from observed mobility.

Early large-scale studies revealed strong regularities in human displacement. Brockmann et al.~\cite{brockmann2006scaling} analyzed the geographic circulation of banknotes in the United States and observed displacement patterns resembling L\'evy flights~\cite{levy_flight,chechkin2008levy}, characterized by many short movements and fewer long-distance displacements. Gonz\'alez et al.~\cite{gonzalez2008understanding}, using mobile-phone records, identified similar heavy-tailed characteristics at the individual level and showed that people repeatedly return to a limited set of locations rather than moving randomly throughout geographic space.

Song et al.~\cite{natureSong2010exploration} formalized this combination of recurrent and novel movement through an exploration and preferential-return model. Individuals may visit new locations, but the probability of exploration decreases as their set of previously visited locations grows, while frequently visited places become increasingly likely destinations. Human mobility is therefore stochastic while remaining strongly structured by recurrent locations.

Mobility regularities also depend on the spatial and temporal scale at which movement is observed. Alessandretti et al.~\cite{natureAlessandretti2017multiScaleSpatioTemporal} demonstrated variations in displacement and waiting-time distributions across different observation scales. In later work, Alessandretti et al.~\cite{alessandretti2020nature} represented mobility through hierarchical spatial ``containers'', associating locations with progressively broader spatial entities. This formulation relates human movement to hierarchical conceptions of geographic space, in which places can be understood through nested structures such as neighborhoods, cities, and regions~\cite{naturePaasi2004placeAndRegion,natureBatty2006hierarchy,naturePumain2006hierarchy}. Scale-free mobility patterns can then emerge from the aggregation of movements occurring across different levels.

These empirical models depend on observations from sources such as mobile-phone records, transportation systems, and location-based services. Such datasets vary in temporal and spatial resolution, population coverage, representativeness, accessibility, and privacy constraints. Haraguchi et al.~\cite{haraguchi2022mobilityDataReview}, in a systematic review of mobility data for urban resilience, emphasize both the increasing availability of mobility information and limitations associated with sampling, coverage, privacy, and transferability. Detailed patterns inferred for one population or region may therefore not be available or directly transferable to another.

There is also a distinction between describing recurring mobility and representing how movement changes when urban conditions are modified. Statistical mobility models can reproduce displacement, visitation, or destination-choice patterns, but exceptional events, restrictions, emergencies, service interruptions, and infrastructure failures may alter the conditions under which those movements occur. In such circumstances, population flows depend not only on established mobility preferences but also on the availability and operation of destinations and services. Mobility data have therefore become relevant to the study of urban resilience under disruptive conditions~\cite{haraguchi2022mobilityDataReview}, while changes in mobility have also been associated with processes such as epidemic propagation~\cite{natureKramer2020mobilityCovid,hernandez2021human}. 

These applications illustrate that mobility can be both an outcome of changing urban conditions and a mechanism through which other processes evolve. They also motivate simulation approaches in which population movement responds to changes in the destinations, services, and conditions represented in the urban environment.

%The COVID-19 pandemic provided a prominent example of this interaction. Kraemer et al.~\cite{natureKramer2020mobilityCovid} related changes in human mobility and control measures to epidemic spread in China. Jamshidi et al.~\cite{jamshidi2020global} examined relationships among mobility, urban density, social isolation, environmental conditions, and infection incidence at different spatial scales, while Hern\'andez-Orallo and Armero-Mart\'inez~\cite{hernandez2021human} used urban mobility simulation to investigate infection risk in spaces such as public squares and transportation facilities. Similar dependencies appear in urban resilience, where disruptions to infrastructure and services may redirect or constrain population flows~\cite{haraguchi2022mobilityDataReview}. These applications show that mobility can be both an outcome of urban conditions and a mechanism through which other processes evolve.

Recent studies increasingly combine conventional mobility formulations with data-driven techniques. Liu et al.~\cite{liu2024GraviGBM} proposed GraviGBM, integrating a gravity-based formulation with gradient-boosted decision trees to predict population flows while preserving elements of interpretability. Guo et al.~\cite{guo2025datadrivendiscoveryhuman} use symbolic regression to derive interpretable mathematical expressions for human mobility directly from data. These approaches seek predictive performance while retaining explicit relationships that can be examined rather than relying exclusively on opaque learned representations.

%Generative methods have also begun to move from mobility prediction toward simulation. CAMS~\cite{du2025camscitygptpoweredagenticframework}, for example, combines large-language-model-based agents with an urban environment to generate mobility patterns from autonomous activity and destination choices. This differs from statistical mobility laws by allowing individual agents to generate activities and trajectories rather than sampling movement exclusively from predefined distributions. The simulated population nevertheless remains organized around individual agents making mobility decisions.

Population mobility research consequently provides strong empirical foundations for large-scale movement, including heavy-tailed displacements, recurrent destinations, exploration, preferential return, and hierarchical spatial organization~\cite{brockmann2006scaling,gonzalez2008understanding,natureSong2010exploration,alessandretti2020nature}. Data-driven methods further extend the ability to infer and predict these patterns~\cite{liu2024GraviGBM,guo2025datadrivendiscoveryhuman}. Simulating urban populations, however, additionally requires mechanisms for generating population activities and for representing the locations, services, and environmental conditions with which those movements interact. These requirements are addressed more directly by the urban population simulation models discussed in the following section.

\section{Urban Population Simulation}
\label{related:multilevel_urban_simulation}

Urban population simulation models represent population activities and movements within spatially structured urban systems. Compared with mobility models that primarily characterize observed movement patterns, these approaches explicitly represent populations, environments, and the processes through which population activities unfold over time. Depending on their objectives, models may operate at different spatial, temporal, and behavioral scales and may represent populations as individuals, groups, or combinations of different levels of detail~\cite{serana2023multilevelmodeling}. Their environments may include geographically distributed destinations, transportation networks, services, and other urban structures through which population activities are organized.

Transportation simulation frameworks provide one established approach to representing urban population movement. MATSim~\cite{balmer2009matsim,balmer2011matsim} represents travelers as individual agents performing daily plans composed of activities and trips through transportation networks. Plans are repeatedly executed and evaluated, allowing alternative travel choices and transportation policies to be examined at large population scales. This individual-level representation preserves detailed information about activities, routes, and travel decisions, while scenario construction requires corresponding descriptions of populations, activity plans, travel demand, and transportation infrastructure.

SUMO~\cite{krajzewicz2012sumo} instead focuses primarily on microscopic road traffic. Vehicles are represented individually and move through explicit road networks according to routing and traffic models, with support for intersections, signals, vehicle types, public transportation, and traffic-control mechanisms. This provides detailed simulation of vehicle and network dynamics, although the central abstraction is traffic flow rather than the broader set of activities and state changes that may characterize an urban population.

General-purpose agent-based simulation environments provide greater flexibility concerning the entities and processes represented. GAMA~\cite{taillandier2019gama}, for example, supports spatially explicit models, GIS information, multiple agent types, and configurable interactions among agents and their environments. Rather than prescribing a particular population representation or mobility process, such platforms provide modeling constructs from which application-specific agents, environments, movement mechanisms, and interactions can be defined. This flexibility supports a broad range of urban applications, although the corresponding population abstractions and domain processes must be specified as part of each model.

Recent generative approaches provide another means of producing urban population activity and mobility. CAMS~\cite{du2025camscitygptpoweredagenticframework}, for example, combines large-language-model-based agents with an urban spatial representation to generate activities, destination choices, and mobility trajectories. Rather than sampling movement exclusively from predefined statistical distributions, individual agents produce mobility behavior from their profiles and information about the surrounding city. This provides a flexible mechanism for generating individual mobility patterns, while retaining the individual agent and its trajectory as the principal population representation.

Urban simulation may also involve multiple representations within the same model. Multilevel modeling provides a general perspective for combining processes or entities represented at different levels of detail~\cite{serana2023multilevelmodeling}. As discussed in Section~\ref{related:crowd_simulation}, crowd models may transition between microscopic and aggregated population representations according to spatial or computational requirements. At an urban scale, similar concerns arise when processes with different spatial or temporal resolutions must coexist. The relevant level of representation may depend on whether the model emphasizes individual travel decisions, collective mobility, local interactions, or population-wide state changes.

Urban population simulation may additionally require the represented environment itself to change during a simulation. Urban-resilience research considers cities as systems composed of interacting socio-ecological and socio-technical components whose ability to maintain or recover their functions may be affected by disruptions~\cite{meerow2016urbanResilience}. Critical-infrastructure research further emphasizes that urban systems do not operate independently: dependencies among power, water, transportation, communication, health care, and other services may allow the disruption of one component to affect the operation of others~\cite{rinaldi2001infrastructureInterdependencies}. Modeling and simulation have consequently been used to examine the behavior and failure propagation of interdependent infrastructure systems~\cite{ouyang2014interdependentInfrastructure}.

Such environmental changes may directly affect population dynamics. Infrastructure failures, service interruptions, capacity limitations, or changes in destination availability can constrain, redirect, or generate population movement~\cite{haraguchi2022mobilityDataReview}. From a population-simulation perspective, urban locations may therefore represent more than geographic destinations. Their semantic roles, capacities, schedules, operating conditions, and relationships with other locations or infrastructure may determine whether population demand can be satisfied and how populations respond when normal conditions change.

The reviewed models consequently emphasize different aspects of urban population simulation. MATSim represents detailed individual activity and transportation plans, SUMO focuses on vehicle-level traffic dynamics, GAMA provides general mechanisms for constructing spatial agent-based models, and CAMS generates individual mobility behavior through language-model-based agents. Multilevel approaches further demonstrate that different population representations can coexist when different levels of detail are required. Research on resilience and infrastructure extends these concerns by showing that the state of the urban environment itself may influence population processes. Together, these approaches establish urban population simulation as a problem involving not only how populations move, but also how they interact with structured and potentially changing urban environments. These distinctions provide the basis for positioning LODUS in the following section.

\section{Positioning LODUS across Crowd, Mobility, and Urban Simulation Models}

The reviewed literature addresses complementary aspects of urban population dynamics. Crowd-simulation models emphasize local movement and interaction, with aggregated and multilevel approaches providing alternatives for representing larger populations~\cite{antonitsch2019bioclouds,xiong2010hybrid,serana2023multilevelmodeling}. Human-mobility research operates over broader spatial and temporal scales and identifies regularities in displacement, recurrent visitation, and movement across spatial scales~\cite{brockmann2006scaling,gonzalez2008understanding,natureSong2010exploration,alessandretti2020nature}. Urban simulation frameworks such as MATSim, SUMO, and GAMA provide explicit representations of transportation networks, geographic environments, or spatially situated agents~\cite{balmer2009matsim,krajzewicz2012sumo,taillandier2019gama}. As discussed in the previous section, resilience and infrastructure studies additionally highlight that the availability and operation of urban locations and services may change under disruptive conditions.

LODUS is positioned across these perspectives by combining an aggregated representation of heterogeneous populations with an explicit representation of the urban locations with which those populations interact. Points of interest may express population demand according to their semantic role, schedules, capacities, and current conditions. A school may request students, an event may temporarily attract population to a destination, and a health-care facility may generate service demand subject to its availability. Population dynamics can therefore be expressed through interactions between population characteristics and the state and function of urban locations.

The population representation complements this environment-oriented organization. Populations are maintained as variable-size groups that can be divided when only a subset participates in a movement or state change and combined again when maintaining their distinction is no longer necessary. This is related to aggregate and multilevel population representations~\cite{antonitsch2019bioclouds,legion,xiong2010hybrid,dang2023multilevelCrowd}, but changes in granularity occur within the population representation rather than through transitions between distinct microscopic and macroscopic simulation models. The same population and environment abstractions can consequently support different population processes without requiring separate underlying representations for each application.

The formulation developed in this thesis builds on the original LODUS framework~\cite{silva2020lodus} and the population-simulation mechanisms subsequently developed as LodusPop~\cite{silva2026lodusPop}. The original framework combined a structured urban environment with BioClouds~\cite{antonitsch2019bioclouds}%SO: acho que vale dizer o que é? bioclouds é um agregado de individuos porém seu movimento é contínuo...
, an aggregated crowd model in which populations are represented as groups of individuals that move continuously through space while sharing collective characteristics such as destination, velocity, and desired density. The subsequent formulation generalized this population representation through heterogeneous variable-size groups whose composition and state can change over time, without requiring continuous spatial movement. 

This thesis consolidates these concepts under the name LODUS and extends their application to dynamic location availability, infrastructure dependencies, and a broader range of urban population processes. The resulting model is also extensible, allowing additional population characteristics, location semantics, behaviors, and domain-specific processes to be incorporated while retaining the same core population and environment abstractions.
%SO: além do extensível

LODUS is therefore complementary to established crowd, mobility, and urban simulation approaches rather than a replacement for them. It does not seek to reproduce detailed pedestrian interactions, vehicle-level traffic dynamics, or individual trajectories inferred from fine-grained mobility data. Instead, it targets an intermediate modeling space in which large populations retain selected heterogeneous characteristics, move among semantically defined urban locations, and interact with environments whose conditions may change over time. This provides the basis for the model presented in the following chapter.